\chapter{Empirical Results} \label{ch:results}

This chapter presents the measured transaction sizes and cost data that form the
empirical basis for the security analysis in Chapter~\ref{ch:security}.
Section~\ref{se:lifecycle_measurements} reports the core lifecycle vsize measurements.
Section~\ref{se:opvault_correction} documents the correction of prior OP\_VAULT
estimates. Section~\ref{se:batching_results} presents batching efficiency data.
Section~\ref{se:revault_results} measures revault amplification costs.
Section~\ref{se:address_reuse_results} reports address reuse behavior.
Section~\ref{se:simplicity_measurements} presents the Simplicity vault measurements
separately.

% ======================================================================
\section{Lifecycle Cost Measurements} \label{se:lifecycle_measurements}

Table~\ref{tab:lifecycle} shows the measured vsize for each transaction in the
deposit--trigger--withdraw lifecycle across all five covenant designs. These
measurements are the foundation for all subsequent cost analysis: every attack cost
function, fee sensitivity projection, and cross-covenant comparison in this thesis
derives from these structural vsize values.

\begin{table}[t]
\centering
\caption{Lifecycle vsize measurements for the four Bitcoin
covenant designs. Each value is regtest-measured and structurally
deterministic (range = 0 across repeated measurements). Lifecycle
total = tovault + trigger + withdraw (the happy-path cost). Simplicity
(Elements) measurements appear in Appendix~C. The \emph{axis measured}
column identifies the primary axis
(Table~\ref{tab:design_space_position}) whose position determines the
row's cost structure.}
\label{tab:lifecycle}
\begin{tabular}{@{}lrrrrcl@{}}
\toprule
Covenant & tovault & trigger & withdraw & recover & \textbf{total}  & axis measured                                                 \\
         & (vB)    & (vB)    & (vB)     & (vB)    & \textbf{(vB)}   &                                                               \\
\midrule
CTV        & 122 &  94 & 152 & 133 & \textbf{368} & \axfee{}=anchor, \axwdbind{}=digest          \\
CCV        & 300 & 154 & 111 & 122 & \textbf{565} & \axrevault{}=yes, \axfee{}=dynamic           \\
OP\_VAULT  & 154 & 292 & 121 & 246 & \textbf{567} & \axrevault{}=yes, \axfee{}=dynamic (wallet)  \\
CAT+CSFS   & 122 & 221 & 210 & 125 & \textbf{553} & \axwdbind{}=dual-bind, \axfee{}=ACP          \\
\bottomrule
\end{tabular}
\end{table}

CTV has the cheapest lifecycle at 368\vB{}, owing to its minimal script structure: the
trigger (unvault) transaction is a bare CTV hash check with no signature, producing
the smallest trigger of any design at 94\vB{}. CCV and OP\_VAULT have nearly identical
lifecycle totals (565 and 567\vB{}), but with opposite cost distributions: CCV's
trigger is cheap (154\vB{}) and its deposit is expensive (300\vB{} due to the
two-output contract instance structure), while OP\_VAULT's trigger is expensive
(292\vB{} due to the fee-wallet input) and its deposit is moderate (154\vB{}).
CAT+CSFS falls between CTV and CCV/OP\_VAULT at 553\vB{}, with the dual CSFS+CHECKSIG
introspection pattern adding witness overhead to both trigger (221\vB{}) and withdrawal
(210\vB{}).

The lifecycle cost at a given fee rate~$r$ (in sat/vB) is:

\begin{equation} \label{eq:lifecycle_cost}
C_{\text{lifecycle}}(r) = (\text{vsize}_{\text{tovault}} + \text{vsize}_{\text{trigger}}
  + \text{vsize}_{\text{withdraw}}) \times r
\end{equation}

\noindent
At the median 2024 fee rate of 50\satvB{}, a single vault lifecycle costs 18{,}400~sats
(CTV), 28{,}250~sats (CCV), 28{,}350~sats (OP\_VAULT), or 27{,}650~sats (CAT+CSFS).

% ======================================================================
\section{OP\_VAULT Vsize Correction} \label{se:opvault_correction}

Prior to this work, the only published transaction size estimates for OP\_VAULT were
informal hand-calculations. Table~\ref{tab:opvault_correction} compares these estimates
against our regtest measurements.

\begin{table}[t]
\centering
\caption{OP\_VAULT vsize corrections. Prior estimates are from informal community
discussion; measured values are from regtest lifecycle runs on the
\texttt{opvault-demo} reference implementation.}
\label{tab:opvault_correction}
\begin{tabular}{@{}lrrr@{}}
\toprule
Transaction & Prior estimate & Measured & Delta \\
\midrule
trigger (\texttt{start\_withdrawal})  & ${\sim}200$\vB{} & 292\vB{} & $+46\%$ \\
recovery (\texttt{OP\_VAULT\_RECOVER}) & ${\sim}170$\vB{} & 246\vB{} & $+45\%$ \\
\bottomrule
\end{tabular}
\end{table}

The discrepancy originates from the fee-input requirement described in
Section~\ref{se:opvault}. OP\_VAULT's value preservation rule prohibits spending the
vault UTXO alone (the fee would violate input--output balance), so every trigger and
recovery transaction includes a second input from a fee wallet. This adds a P2TR
input witness (approximately 80--90\vB{}) that prior estimates omitted. The
consequence is a 36\% lifecycle cost premium over CCV: OP\_VAULT's lifecycle
(567\vB{}) versus CCV's (565\vB{}) appears similar in total, but OP\_VAULT pays
80--90\vB{} more for every non-deposit transaction. Over time, the accumulated
overhead is significant for vault operators who trigger frequently.

This measurement contributed to the economic case for withdrawing BIP-345 in favor
of BIP-443~\cite{OB25}: the purpose-built opcodes' fee-input overhead costs more
than the generic CCV approach achieves with the same vault semantics.

% ======================================================================
\section{Batching Efficiency} \label{se:batching_results}

CCV and OP\_VAULT can combine multiple vault UTXOs into a single trigger transaction.
CTV and CAT+CSFS cannot batch because each vault's template hash (CTV) or embedded
\texttt{sha\_single\_output} (CAT+CSFS) commits to a specific output structure that
changes with each additional input.

Table~\ref{tab:batching} shows the measured trigger vsize at increasing batch sizes.
Both CCV and OP\_VAULT exhibit linear scaling with a fixed base cost and a constant
marginal cost per additional vault.

\begin{table}[t]
\centering
\caption{Batched trigger vsize measurements. CCV and OP\_VAULT are measured from
actual multi-input transactions on regtest. CTV and CAT+CSFS cannot batch.}
\label{tab:batching}
\begin{tabular}{@{}lrrrrrr@{}}
\toprule
& \multicolumn{6}{c}{Batch size ($N$ vaults)} \\
\cmidrule(l){2-7}
Covenant & $N{=}1$ & $N{=}2$ & $N{=}3$ & $N{=}5$ & $N{=}10$ & $N{=}20$ \\
\midrule
CCV       & 154 & 254  & 355   & 555   & 1{,}056  & 2{,}059 \\
OP\_VAULT & 292 & 385  & 479   & 667   & 1{,}135  & 2{,}073 \\
\midrule
\multicolumn{7}{@{}l}{\footnotesize CTV, CAT+CSFS: $N \times$ single trigger
(154\vB{}, 221\vB{} respectively)} \\
\bottomrule
\end{tabular}
\end{table}

The batching efficiency follows a linear model:

\begin{equation} \label{eq:batching}
\text{vsize}_{\text{batch}}(N) = \text{base} + \text{marginal} \times N
\end{equation}

\noindent
For CCV, the marginal cost is ${\sim}100$\vB{} per additional vault (base
${\sim}54$\vB{}). For OP\_VAULT, the marginal cost is ${\sim}94$\vB{} per vault
(base ${\sim}198$\vB{}). The marginal cost per vault in a batch is substantially
lower than the cost of individual triggers (154\vB{} for CCV, 292\vB{} for
OP\_VAULT), demonstrating the efficiency gain from amortizing the transaction
overhead across multiple inputs.

The standardness ceiling---the maximum batch size that fits within Bitcoin Core's
400{,}000~WU transaction weight limit---can be projected from these measurements:

\begin{equation} \label{eq:ceiling}
N_{\max} = \left\lfloor \frac{400{,}000 - \text{base\_weight}}
  {\text{marginal\_weight}} \right\rfloor
\end{equation}

\noindent
For CCV, $N_{\max} \approx 1{,}000$ vaults per transaction. For OP\_VAULT,
$N_{\max} \approx 1{,}064$ vaults. In practice, batch sizes above 50--100 are rare
in custodial operations, so the standardness limit does not constrain realistic
deployments.

% ======================================================================
\section{Revault Amplification} \label{se:revault_results}

CCV and OP\_VAULT support partial withdrawals: the vault owner withdraws a portion
of the vault balance and re-locks the remainder in a new vault UTXO. This process
can be repeated, creating a chain of partial withdrawals from a single initial
deposit.

The cost of each withdrawal step is the sum of the trigger (or trigger-and-revault)
and the withdrawal completion. Table~\ref{tab:revault} shows the per-step and
cumulative costs measured across 9~successive partial withdrawals of 5{,}000{,}000~sats
each from an initial 49{,}999{,}900-sat vault.

\begin{table}[t]
\centering
\caption{Revault amplification. Per-step withdrawal vsize is constant across steps
because the script structure does not change with balance. CTV and CAT+CSFS cannot
revault natively; their cost is the full unvault--withdraw cycle per withdrawal.}
\label{tab:revault}
\begin{tabular}{@{}lrrr@{}}
\toprule
Covenant & Per-step vsize & Cumulative (5 steps) & Cumulative (9 steps) \\
\midrule
CCV       & 111\vB{} & 555\vB{} & 999\vB{} \\
OP\_VAULT & 121\vB{} & 605\vB{} & 1{,}089\vB{} \\
\midrule
CTV (extrapolated)      & 246\vB{} & 1{,}230\vB{} & 2{,}214\vB{} \\
CAT+CSFS (extrapolated) & 431\vB{} & 2{,}155\vB{} & 3{,}879\vB{} \\
\bottomrule
\end{tabular}
\end{table}

The per-step vsize is constant across all 9~steps for both CCV and OP\_VAULT, confirming
the structural determinism of the measurement. CCV's smaller per-step cost (111\vB{}
vs.\ 121\vB{}) reflects its lighter withdrawal transaction (no fee-input overhead).

For CTV and CAT+CSFS, which lack native revault, each partial withdrawal requires a
full cycle: unvault the entire balance, withdraw the desired portion, and re-deposit
the remainder in a new vault. The extrapolated per-step cost is the sum of the trigger
and withdrawal vsizes for a single cycle. This is 2.2$\times$ (CTV) to 3.9$\times$
(CAT+CSFS) more expensive than CCV's native revault per step.

The cumulative cost has implications for the per-UTXO recovery-cost scaling attack analyzed in
Chapter~\ref{ch:security}: an attacker chaining many small partial withdrawals forces
the watchtower to pay recovery fees proportional to the number of splits.

% ======================================================================
\section{Address Reuse} \label{se:address_reuse_results}

CTV is the only design where address reuse causes permanent fund loss. A second deposit
to the same CTV address creates a UTXO whose amount does not match the committed
template hash. The template commits to output amounts derived from the \emph{original}
deposit, so the second UTXO is either permanently unspendable (if the deposit is
smaller) or overpays miners by the difference (if larger).

CCV, OP\_VAULT, CAT+CSFS, and Simplicity all handle multiple deposits to the same
address safely. CCV and OP\_VAULT use per-instance contract checking---each UTXO is
independently validated against the covenant rules regardless of how many UTXOs share
the same address. CAT+CSFS generates a unique vault plan per source coin, producing
distinct P2TR addresses. Simplicity creates fresh plans with new addresses for each
vault instance.

This is a design-level property, not an implementation artifact: CTV's template hash
commits to input count and output amounts, making it structurally incompatible with
variable deposit amounts. The other designs either validate per-input (CCV, OP\_VAULT)
or per-instance (CAT+CSFS, Simplicity).

% ======================================================================
\section{Fee Sensitivity Projections} \label{se:fee_sensitivity_results}

Experiment~J projects the measured vsize data from the preceding sections into
real-world fee environments using the linear fee model
(Equation~\ref{eq:fee_projection}). Table~\ref{tab:fee_sensitivity} shows the
lifecycle cost and key attack costs across six historically observed fee rates.

\begin{table}[t]
\centering
\caption{Fee sensitivity projections. Lifecycle cost =
$\text{vsize}_{\text{total}} \times r$. Fee pinning cost is fee-invariant.
Recovery griefing cost is the defender's per-round cost. Watchtower exhaustion
shows the number of splits needed to exhaust a 0.5~BTC vault to dust.}
\label{tab:fee_sensitivity}
\footnotesize
\begin{tabular}{@{}rrrrrrrr@{}}
\toprule
& \multicolumn{4}{c}{Lifecycle cost (sats)} & Fee pin & Grief & WT exhaust \\
\cmidrule(lr){2-5}
$r$ (sat/vB) & CTV & CCV & OPV & CAT & (CTV) & (CCV) & splits \\
\midrule
1   &     368 &     565 &     567 &     553 & 14{,}300 &    454 & 409{,}836 \\
10  &   3{,}680 &   5{,}650 &   5{,}670 &   5{,}530 & 14{,}300 &  4{,}540 &  40{,}984 \\
50  &  18{,}400 &  28{,}250 &  28{,}350 &  27{,}650 & 14{,}300 & 22{,}700 &   8{,}197 \\
100 &  36{,}800 &  56{,}500 &  56{,}700 &  55{,}300 & 14{,}300 & 45{,}400 &   4{,}098 \\
300 & 110{,}400 & 169{,}500 & 170{,}100 & 165{,}900 & 14{,}300 & 136{,}200 &  1{,}366 \\
500 & 184{,}000 & 282{,}500 & 283{,}500 & 276{,}500 & 14{,}300 & 227{,}000 &    820 \\
\bottomrule
\end{tabular}
\end{table}

Three observations emerge from the projections:

\begin{itemize}
\item \textbf{Fee pinning is fee-invariant.} The 14{,}300-sat pinning cost
  (Section~\ref{se:fee_pinning}) is constant across all fee environments because
  the attacker pays dust-rate fees regardless of prevailing rates. At a 0.5~BTC
  vault value, pinning costs 0.03\% of the vault---well below any rationality
  threshold at any historical fee rate.

\item \textbf{Lifecycle cost rankings are stable.} CTV is the cheapest at every fee
  rate. CCV and OP\_VAULT are nearly identical (within 0.4\%). CAT+CSFS falls between
  CTV and CCV/OP\_VAULT. The ranking does not change with fees because the vsize
  ratios are constant.

\item \textbf{Watchtower exhaustion feasibility is fee-dependent.} At 1\satvB{},
  exhausting a 0.5~BTC vault requires ${\sim}410{,}000$ splits (infeasible---66 blocks
  of maximum-rate splitting). At 300\satvB{}, it requires only 1{,}366 splits
  (feasible in a single block). This fee-dependent transition is the empirical basis
  for Proposition~\ref{thm:fee_inversion} in Chapter~\ref{ch:security}.
\end{itemize}

% ======================================================================
\section{Synthesis: What the Measurements Imply} \label{se:measurement_synthesis}

Three patterns in the data foreshadow the security analysis in
Chapter~\ref{ch:security}:

\begin{enumerate}
\item \textbf{CTV is the cheapest lifecycle but the least flexible.}
  At 368\vB{}, CTV's lifecycle cost is 33--35\% lower than CCV or OP\_VAULT.
  But CTV cannot batch, cannot partially withdraw, and loses funds on address
  reuse. The cost advantage comes from CTV's minimal script structure---the same
  rigidity that creates its security vulnerabilities (fee pinning via anchors,
  stuck funds on reuse).

\item \textbf{OP\_VAULT's fee-input overhead is structural, not incidental.}
  The 80--90\vB{} per non-deposit transaction reflects BIP-345's value preservation
  rule, not an implementation choice. This overhead makes OP\_VAULT 36\% more
  expensive than CCV per lifecycle---the economic basis for BIP-345's withdrawal
  in favor of BIP-443.

\item \textbf{Batching and revault create the per-UTXO recovery-cost scaling surface.}
  CCV and OP\_VAULT's capability advantages (partial withdrawal, batched triggers)
  are also the attack surface for per-UTXO recovery-cost scaling. The per-step revault cost
  (111--121\vB{}) is cheap enough to chain thousands of splits, creating the
  fee-dependent vulnerability analyzed in Section~\ref{se:watchtower_exhaustion}.
  CTV and CAT+CSFS are immune precisely because they lack these capabilities.
\end{enumerate}

% ======================================================================
\section{Case Study: Simplicity on Elements} \label{se:simplicity_measurements}

The preceding sections present the four-way Bitcoin covenant comparison---the core
contribution of this thesis. This section reports measurements from a fifth vault
built with the Simplicity language on Elements regtest. Because Elements uses
different transaction serialization (explicit fee outputs, confidential transaction
fields, multi-asset model) and a federated consensus model, these measurements are
\textbf{not directly comparable} with the Bitcoin results. We present them as a
case study in what covenant enforcement looks like with full transaction introspection
and no soft-fork dependency.

\paragraph{Lifecycle costs.}
The Simplicity vault lifecycle totals 865\vB{} (tovault 269, trigger 298, withdraw
298, recover 286). This is 49--135\% larger than the Bitcoin designs, reflecting
two factors: Elements serialization overhead (approximately 30--40~bytes per output
for confidential transaction fields even in non-confidential mode) and Simplicity's
bit-oriented DAG witness format (structurally different from Bitcoin Script's
stack-based witnesses). A hypothetical Bitcoin-native Simplicity implementation
would produce smaller transactions; the 865\vB{} figure includes both the language
overhead and the chain overhead.

\paragraph{Output-constrained recovery.}
Unlike CAT+CSFS, the Simplicity vault enforces \texttt{jet::outputs\_hash()} on
\emph{both} trigger and recovery paths. Cold key compromise can only recover funds
to the pre-committed recovery address---matching CCV and OP\_VAULT's output-constrained
model and providing stronger recovery safety than CAT+CSFS's unconstrained
\texttt{OP\_CHECKSIG} (TM11, Section~\ref{se:cold_key_recovery}).

\paragraph{Fee pinning immunity.}
Elements uses explicit fee outputs (empty scriptPubKey) rather than Bitcoin's implicit
fees. There are no anchor outputs to pin with descendant chains, eliminating the
TM1 attack vector entirely. This immunity is a property of the Elements transaction
format, not of Simplicity's covenant mechanism.
